\subsection{An explicit CM repair and the actual adelic endpoint}

\begin{proposition}[DP57: unramified CM repair and endpoint realization]\label{DP57}
Let $E=\mathbf Q(\zeta _7)$, $F=E^+$, and let $J$ be the principally
polarized genus-three Jacobian of DP22 and DP49, of type $\{1,2,3\}$.
Put
\[
 E'=E(i)=\mathbf Q(\zeta _{28}),\qquad F'=(E')^+=F(\sqrt7).
\]
\begin{enumerate}[label=(\alph*),leftmargin=*]
\item The original $E$ fails the literal $\operatorname{IsCMUnram}$
condition used in the source's $\operatorname{CMData}$.
The pair $(F',E')$ satisfies $\operatorname{CMData}$ directly.
The extension $\mathcal O_{E'}/\mathcal O_{F'}$ is finite
\emph{étale}, and the field-level toroidal $\operatorname{Host}$
can also be supplied with $d=7$.
\item The product $A=J^2$ admits an integral $\mathcal O_{E'}$-action
whose Rosati involution for the product principal polarization is
canonical conjugation. Its CM type is
\[
 \Phi'=\{a\in(\mathbf Z/28)^\times:a\bmod7\in\{1,2,3\}\}
      =\{1,3,9,15,17,23\}.
\]
At $11$, its two $E'$-prime factors have heights $6,6$,
dimensions $4,2$, and slopes $2/3,1/3$.
The underlying product retains its height-three summands, but these
are not individual $E'$-stable factors.
\item For every positive Hermitian rank and dominant coefficient
weights, the generic source construction supplies an actual endpoint
level with $T=1$ and principal reduction level at
\[
 P=(29,\zeta _{28}-2),\qquad
 S_T=\varnothing,\quad S_G=\{P\}.
\]
For this fresh level, $\operatorname{DecoPtSurj}$ and
$\operatorname{IsLevOf}$ hold.
The source's finite coefficient-value space is exactly the restriction
of its adelic finite-weight space. It contains a nonzero vector with
an actual detecting partner period equal to $1$.
\item The type $\Phi'$ is induced from no imaginary quadratic
subfield. Thus this construction does not supply the old
$\operatorname{Datum.dom}$ or $\operatorname{CMObj.HasType}$
theta input for this geometric $A$.
The separate old wrapper condition $8\mid f$ also fails, since $f=6$;
neither restriction occurs in the generic endpoint construction in (c).
\end{enumerate}
The construction doubles the geometric dimension and changes the field
and level. Its product polarization has degree one.
No preservation of a prescribed earlier level, quantitative global-cost
bound, or membership in the prescribed iterated theta image is asserted.
\end{proposition}

\begin{proof}
The source's $\operatorname{CMData}$ is
\path{Formalism/LocalRingCardinality/DataSpaces.lean}, 14--19.
Its $\operatorname{IsCMUnram}$ field is explicitly
\[
 \begin{gathered}
 E\text{ totally complex, with conjugation compatible with embeddings},
 \\
 \forall Q\subset\mathcal O_E\text{ maximal},\
       \exists\theta\in\mathcal O_E:
                    \theta-\bar\theta\notin Q.
 \end{gathered}
\]
See \path{Formalism/HermitianLattices/StatementMass.lean}, 44--48.
In $\mathcal O_E=\mathbf Z[\zeta _7]$, the prime
$\mathfrak p=(1-\zeta _7)$ has residue field $\mathbf F_7$.
Conjugation is the identity modulo $\mathfrak p$, so every
$\theta-\bar\theta$ belongs to it. This proves the first failure.

For $E'$ use the explicit norm identity
\[
 N_{E/\mathbf Q}(\zeta _7-\zeta _7^{-1})
 =N(\zeta _7^{-1})N(\zeta _7-1)N(\zeta _7+1)
 =1\cdot7\cdot1=7.
\]
Thus $\theta=\zeta _7$ satisfies the required condition at every
prime of $E'$ not over $7$; at primes over $7$ use $\theta=i$,
whose difference from its conjugate is the unit $2i$.
The field is cyclotomic and totally complex, so the other
conjugation and real-fixed-field conditions hold.
For the exact cyclotomic integer-ring and prime-power assertions,
see J.\ S.\ Milne, \emph{Algebraic Number Theory}, version 3.08,
Proposition 6.2(b),(c), printed pp.~96--98,
and Theorem 6.4(a),(b), printed p.~99,
\url{https://www.jmilne.org/math/CourseNotes/ANT.pdf}.

Write
\[
 z=\zeta _7+\zeta _7^2+\zeta _7^4,\qquad
 z+\bar z=-1,\qquad z\bar z=2,\qquad
 \delta=1+2z,\quad\delta^2=-7.
\]
Then $i\delta$ is real with square $7$.
Since $[F:\mathbf Q]=3$, the field $F(\sqrt7)$ has degree six
and is the fixed field of $E'$.

The same witnesses prove literal unramifiedness of the integer rings.
At a finite prime $v$ of $F'$, let $R=(\mathcal O_{F'})_v$ and
choose $\theta=\zeta _7$ if $v\nmid7$, and $\theta=i$ otherwise.
It generates the quadratic extension, its trace and norm belong to $R$,
and
\[
 \operatorname{disc}(1,\theta)=(\theta-\bar\theta)^2\in R^\times.
\]
Hence $R[\theta]$ is finite étale over $R$ and is normal.
For the latter implication one may use Stacks Project,
Lemma 10.163.9, Tag 033C,
\url{https://stacks.math.columbia.edu/tag/033C}.
Its fraction field is $E'$, and every element integral over $R$
is integral over $R[\theta]$; normality therefore makes this order the
full integral closure. This proves étaleness at every finite prime.
In
\path{ReciprocityTheory/ReductiveGroups/ArithmeticGeometry/BirationalGeometry/ToroidalCompactification/TwoWeightEqOneAddRoot.lean},
14--26, take
\[
 d=7,\qquad rt=-\delta,\qquad w_0=-z,\qquad
 rt_0=\sigma_0(rt),
\]
with canonical conjugation. Then $rt^2=-7$ and $2w_0=1+rt$,
and the just-proved étaleness supplies $\operatorname{Host.unram}$.
This completes the field-level construction in (a).

For (b), let $E$ act diagonally on $A$ and let
\[
 I=\begin{pmatrix}0&-1\\1&0\end{pmatrix}
\]
act on the two factors. It commutes with $E$ and has square $-1$.
The unique quadratic subfield of $E$ is $\mathbf Q(\sqrt{-7})$,
so $i\notin E$ and this gives an injective action of $E(i)$.
With compatible roots,
\[
 \zeta _7=\zeta _{28}^4,\quad i=\zeta _{28}^7,\quad
 \zeta _{28}=\zeta _7^2i^{-1},\qquad
 \mathcal O_{E'}=\mathbf Z[\zeta _7,i].
\]
Both generators act integrally, so this is the full integer-ring
action without any isogeny.
Rosati for the product polarization sends $I$ to
$I^{\mathsf T}=-I$ and $a\in E$ to $\bar a$.
It is therefore canonical conjugation on $E'$.
Equivalently $-\delta I$ is self-adjoint and has square $7$.
The polarization remains principal; the degree $7^6$ arising from
the alternative presentation with polarization
$\operatorname{diag}(1,7)$ is not a necessary cost.

On each old tangent-character multiplicity space, $I$ has the
two distinct eigenvalues $i,-i$. Thus each old selected embedding
has both extensions selected, giving the stated $\Phi'$.
The Frobenius orbits at $11$, with their selected indicators, are
\[
 \begin{array}{c|c}
 (1,11,9,15,25,23)&(1,0,1,1,0,1)\\
 (3,5,27,17,19,13)&(1,0,0,1,0,0).
 \end{array}
\]
The dimensions are four and two; the underlying factors are the
squares of the two height-three factors of $J$ from DP22 and DP49.
This proves their stated heights and slopes.
Each local $E'$-field at $11$ is unramified of degree six.
After splitting its action on a rational Dieudonné module,
Frobenius cyclically identifies the six eigenspaces. They have
equal dimensions, so any nonzero factor stable under this field
has height divisible by six. This rules out treating the old
height-three product summands as individual $E'$-stable factors.

For (c), $2\in\mathbf F_{29}^\times$ has order $28$:
$2^{14}=-1$ and $2^4=16\ne1$.
Thus evaluation $\zeta _{28}\mapsto2$ gives the displayed prime
with residue field $\mathbf F_{29}$.
All primitive $28$-th roots are simple modulo $29$, so its
completion is $\mathbf Q_{29}$.
Conjugation replaces $2$ by $2^{-1}=15$, proving $P\ne\bar P$.
The definitions \path{IsConjPl} and \path{IsSplitSet} in
\path{Formalism/LocalToGlobal/Divisibility/DefGroupLemmas.lean},
23--27, therefore give
$\operatorname{IsSplitSet}(S_T\cup S_G)$.
Disjointness is automatic since $S_T$ is empty.

Hereafter the source paths beginning with \path{Columns/} are under
\path{Formalism/LocalToGlobal/}.
The theorem \path{Columns/InputsGenus.lean},
\path{unimodGenus_holds}, 17--19, supplies
$\operatorname{UnimodGenus}(hc,n)$ for the CM datum just constructed
and every $n>0$.
The theorem \path{reduction_surjective_bis} in the preceding
\path{DefGroupLemmas.lean}, 1475--1481, supplies
$\operatorname{RedSurj}$.
Its definition at 64--68 allows arbitrary invertible reductions
at an $\operatorname{IsSplitSet}$ for a framed definite unimodular lattice.
The theorem \path{Columns/AdelicLemmas.lean},
\path{decoPtSurj_of}, 1245, consumes exactly these two statements,
the split-set condition and disjointness.
It proves the actual $\operatorname{DecoPtSurj}$ input.

Take the finite reduction subgroup $H_P=\{1\}$ with trivial
character. In the generic constructor
\path{Columns/RecLevelStep.lean}, \path{stepOf}, 20--51,
set
\[
 T=1,\qquad \pi=1,\qquad f=6,\qquad S_T=\varnothing.
\]
Divisibility by $\pi$ is automatic, $\pi\bar\pi=T$ is $1=1$,
and $1$ is in no prime. The conditions at the empty set and at
primes dividing $T$ are vacuous.
The standard lattice is self-dual, so its required relative index
is $1=T^{2f}$.
These are precisely the inputs already used by
\path{Columns/TiedXp.lean}, \path{xp_basePi_one}, 21--23,
and \path{tieLevP}, 35--41; the generic constructor accepts the
present arbitrary singleton $S_G$ without an
$\operatorname{Fields.Ok}$ argument.
Its fields are literally
\[
 L.U=\operatorname{levU}(hc,n,\varnothing,S_G,\Lambda.\operatorname{grp}),
 \qquad
 L.\chi=\operatorname{levTyp}(\sigma_0\circ\Lambda.\operatorname{chr}),
\]
at \path{RecLevelStep.lean}, 46 and 50.
These are exactly the two fields of
$\operatorname{IsLevOf}$ at \path{Columns/TiedValFi.lean}, 151--154.
Both equalities follow by definition, just as in that file's
\path{tieLev_isLevOf}, 265--266.
The resulting level is compact and open, principal modulo $29$
at the selected split place and standard hyperspecial elsewhere.

Choose $K=E'$ and the embeddings $\Phi'$ as the coefficient data.
They satisfy the conjugation compatibilities of DP52.
Its weak-goodness proof supplies
$\operatorname{formGSD.Ok}$, with no
$\operatorname{Fields.Ok.quad}$ requirement.
All hypotheses of \path{Columns/TiedValFi.lean},
\path{fi_iff_FiC}, 239--255, are now supplied:
goodness, $\operatorname{DecoPtSurj}$ and
$\operatorname{IsLevOf}$.
That theorem identifies $\operatorname{Fi}$ with the restrictions
of actual adelic finite-weight forms.
Its forward direction uses \path{adOf}, 179--181, and
\path{adOf_mem}, 185--201; its reverse direction uses
the precise stabilizer-invariance theorem at 209--237.

At the chosen prime the principal congruence group is
$1+29M_n(\mathbf Z_{29})$, which is torsion-free by DP55.
Every decorated positive-definite lattice has finite stabilizer,
so the principal decoration makes that stabilizer trivial.
The standard decorated lattice and the nonzero extremal coefficient
give a nonzero finite value supported on a free class.
It is the restriction of an actual adelic form by the preceding
identification. DP55's perfect partner pairing and period formula
then produce a partner with period $1$.
This proves (c) at the fresh level constructed here, rather than
asserting preservation of some earlier fixed level.

Finally, $3,15\in\Phi'$ and
\[
 \langle3\rangle=\{1,3,9,27,25,19\},\qquad
 15\notin\langle3\rangle,\qquad15^2=1.
\]
They generate the full twelve-element Galois group.
If $\Phi'$ restricted to one embedding of an imaginary quadratic
subfield, these automorphisms, compared with the identity in
$\Phi'$, would fix that subfield. It would therefore be contained
in $\mathbf Q$, a contradiction. The common restriction field
$K_q$ is consequently $\mathbf Q$.
The source's uniform two-type condition is
\path{Development/AbelianVarietyTheory/Definitions/CMLattices/IsDiffOf.lean},
\path{HasType}, 64--67, with
$\operatorname{CMObj.typeP,typeM,typeB}$ at 115--122.
Its matching trace condition is $\operatorname{Datum.dom}$ at
\path{TwoWeightEqOneAddRoot.lean}, 43--49.
The elliptic Serre tensor tie $\operatorname{IsTen}$ at
\path{IsDiffOf.lean}, 73--77, is not furnished by an adelic
endpoint identification.
The separate $8\mid f$ requirement is
\path{Columns/AbelianHostFieldConditions.lean}, 110--120,
and fails for $f=6$.
These facts establish (d).
\end{proof}

The unramified CM datum, the field-level toroidal host, the actual
adelic endpoint and a nonzero detecting period have all been constructed.
The remaining mixed-type input is the particular theta family and its
coefficient, path and nonvanishing identities, with the corresponding
global costs. The doubled geometry and its old height-three product
factors do not furnish those identities or a shorter lift chain by
themselves.
